\section{Conclusion}
\label{sec:conclusion}

We presented a systematic multi-level evaluation of MDL-to-UsdPreviewSurface
material simplification in NVIDIA Isaac Sim, quantifying trade-offs across
three complementary dimensions: pixel-level image quality, rendering performance,
and feature-level semantic preservation.

Our pixel-level analysis demonstrates that UsdPreviewSurface conversions
achieve high fidelity relative to the original MDL renderings, with a mean
PSNR of 35.52\,dB, SSIM of 0.990, and LPIPS of 0.020 across 24 matched image
pairs. These results indicate that the four primary material channels ---
diffuse albedo, roughness, metallic, and normal --- capture the dominant visual
appearance of the tested assets, with degradation concentrated in scenes
featuring complex multi-layer MDL effects such as procedural noise textures and
specular coatings.

The performance outcome depends strongly on workload. In the single-object
furniture benchmark, rendering throughput and GPU memory are nearly
indistinguishable between MDL and UsdPreviewSurface, with only reduced
first-load latency variability for the converted assets. In a supplemental
whole-scene benchmark on a large GRScenes commercial interior, however, noMDL
reduces time to a fully warmed ready-to-render state from
$52.02 \pm 2.69$\,s to $13.70 \pm 1.01$\,s and lowers post-load GPU memory by
about 199\,MB. The practical conclusion is therefore two-part: conversion does
not buy much steady-state rendering throughput for shallow single-object scenes,
but it can substantially improve startup behavior once scene scale and material
complexity are high.

Feature-level evaluation confirms that semantic content is well preserved:
CLIP embeddings achieve a mean cosine similarity of 0.925 across matched pairs,
while DINOv2 --- which is more sensitive to fine-grained texture differences ---
yields 0.872. These higher-level metrics are particularly important for
practitioners generating synthetic training data, as they directly probe
whether vision models perceive the simplified renderings as equivalent to the
originals.

The GRScenes pilot extends this analysis to language-conditioned grounding. On
a 15-pair clean visual-QA pool, Gemma4 preserves answer accuracy across all
original and converted renders, while point grounding remains less stable. On a
frozen 30-pair target-centered material-shift stress set, Gemma4 is stable under
the normalized-1000 protocol, while Qwen2.5-VL shows answer and raw-point
degradation on converted renders and stronger raw-coordinate than
normalized-coordinate behavior. These results support the use of VLM probes as
diagnostics for material conversion and as a controlled stress set, but not yet
as broad downstream robustness evidence.

Based on our findings, we offer the following preliminary guidelines for
synthetic data practitioners:
\begin{itemize}
    \item For the tested assets with standard PBR materials (diffuse,
    roughness, metallic, normal maps), MDL-to-UsdPreviewSurface conversion
    appears low-risk under the tested visual and feature metrics
    (PSNR $>$ 35\,dB, SSIM $>$ 0.99).
    \item Assets with complex multi-layer MDL effects should be validated
    individually, as conversion quality varies with material complexity (PSNR
    range: 23--44\,dB across viewpoints).
    \item The conversion enables broader toolchain compatibility --- including
    standard USD viewers, web-based visualization, and GLB export --- without
    sacrificing visual quality for most use cases.
    \item For VLM or embodied-language use cases, clean visual-QA pools and
    material-shift stress pools should be reported separately before claiming
    benchmark-level robustness.
\end{itemize}

\paragraph{Limitations.}
Our evaluation is limited to four chest-of-drawers assets from the Isaac Sim
library, all belonging to the same furniture category. While these assets span
a range of material complexity, further validation on diverse object categories
(e.g., transparent or translucent materials, highly specular metals, organic
surfaces) is needed to establish broader generalizability. Additionally, our
performance benchmarks use single-object scenes; the efficiency impact of
material simplification beyond the one GRScenes commercial-interior startup
benchmark remains an open question. Our performance claims are therefore
limited to four single-object furniture scenes plus one GRScenes commercial
interior under a headless RTX path-tracing setup, and should not yet be
generalized to all larger, heavily instanced, or materially diverse
environments. The VLM results are also claim-bounded: the clean pool remains
below the clean final gate, the stress set is target-centered and category
imbalanced, and Qwen coordinate semantics are unresolved.

\paragraph{Future Work.}
We plan to extend this evaluation along several axes: (1)~scaling to diverse
object categories and material types; (2)~scaling the GRScenes VLM clean pool
to a final preservation benchmark and adding downstream navigation or
manipulation proxies where appropriate; (3)~investigating the cumulative effect
of material simplification combined with mesh simplification and texture
compression in the full USD-to-GLB pipeline; and
(4)~benchmarking rendering in complex multi-object scenes where material count
may become a bottleneck.
